\section{Use the magnitude of the second Pell coordinate}
\label{sec:pell}

\subsection{The precise target}

In \code{NativePellBootstrap}, the explicit solution
\[
  a=(2Q+1,2),\qquad d=Q(Q+1),\qquad Q\in\N,\ Q>0
\]
is proved fundamental for the Pell equation
\[
  x^2-dy^2=1.
\]
The statement \code{firstUnit_fundamental} is used by
\code{first_norm_classification}; the latter turns actual norm solutions
into powers of the fundamental unit \cite{pellsource}.

The existing proof establishes \(y^2\ge4\) for a competing solution with
\(x>1\). Assuming \(y^2<4\), it derives the alternatives
\(y=-1\), \(y=0\), and \(y=1\), then handles them separately. The two
nonzero cases have the same norm equation. The zero case is already ruled
out by a theorem in Mathlib's Pell API.

\subsection{The shorter argument}

For a competing solution \(b=(x,y)\) with \(x>1\), use
\code{Pell.Solution₁.y_ne_zero_of_one_lt_x}. If \(y^2<4\), this gives
\[
  0<|y|<2.
\]
Integer discreteness immediately gives \(|y|=1\), hence \(y^2=1\).
No sign of \(y\) has to be selected. Substitution into the norm equation
yields
\[
  x^2=Q(Q+1)+1=Q^2+Q+1.
\]
Because \(Q>0\),
\[
  Q^2<x^2<(Q+1)^2.
\]
Since \(x>1\), comparison of nonnegative squares implies
\(Q<x<Q+1\), impossible for integers. Therefore \(4\le y^2\).
The norm identity now gives
\[
  (2Q+1)^2=1+4Q(Q+1)\le 1+Q(Q+1)y^2=x^2,
\]
so \(2Q+1\le x\). Together with the positivity of the displayed solution,
this is exactly the fundamental-unit condition.

\Needspace{9\baselineskip}
The central replacement fragment is:
\begin{lean}
  have hy0 := abs_pos.mpr
    (Pell.Solution₁.y_ne_zero_of_one_lt_x hb)
  have hylt := abs_lt_of_sq_lt_sq
    (by omega : b.y ^ 2 < (2 : Int) ^ 2)
    (by norm_num : (0 : Int) ≤ 2)
  have hy : |b.y| = 1 := by omega
  have hsq : b.y ^ 2 = 1 := by nlinarith [sq_abs b.y]
\end{lean}

The supplied production patch adds no helper or import. Only the proof of
\code{firstUnit_fundamental} changes. Its public statement and the
\code{firstUnit} definition are unchanged. The full file decreases from
671 to 663 physical lines, with normal blank separation retained.

\subsection{Optional general criterion}

The archive also supplies a checked architectural prototype:
\begin{lean}
theorem fundamental_of_minimal_y_sq
    (a : Pell.Solution₁ d)
    (hax : 1 < a.x) (hay : 0 < a.y)
    (hmin : ∀ b : Pell.Solution₁ d,
      1 < b.x → a.y ^ 2 ≤ b.y ^ 2) :
    Pell.IsFundamental a
\end{lean}

Its proof transfers the inequality on second-coordinate squares through
\code{a.prop_x} and \code{b.prop_x}, then compares positive first
coordinates. The positivity of \(d\) follows from the existing positive
nontrivial Pell solution. The hypothesis deliberately concerns \(y^2\)
for all competing solutions with \(x>1\); it does not silently assume that
the competitor's \(y\) is positive.

For this single client, the direct patch has the better immediate cost.
The generic criterion belongs in a shared module only after another client
justifies its abstraction cost. The optional prototype's helper lines are
not included in the claimed eight-line production reduction. Its separate
square-gap helper is explanatory; a future integration should reuse or
relocate the already existing private square-gap result in
\code{MinimumNorm}, rather than maintain two equivalent long-term helpers
\cite{minimumnorm}.

\subsection{Novelty boundary}

Prior Pell proposals address parity, congruence up to sign, source-register
facts, and rational-root descent. This proposal concerns the different
minimality obligation defining the fundamental solution. It still applies
after those other improvements. The mathematical fact about consecutive
squares is elementary; the new contribution here is the source-specific
removal of the signed enumeration and identification of the smaller proof
obligation.
